% !TEX root = ../../ctfp-print.tex

\lettrine[lhang=0.17]{我}{们已讨论过}函子(functor)作为保持其结构的范畴(category)间映射。

一个函子会将一个范畴"嵌入"另一个范畴中。它可能将多个事物合并为一，但从不会破坏连接关系。一种理解方式是：通过函子我们正在一个范畴内部对另一个范畴进行建模。源范畴充当了一个模型，作为目标范畴中某个结构的蓝图(blueprint)。

\begin{figure}[H]
  \centering\includegraphics[width=0.4\textwidth]{images/1_functors.jpg}
\end{figure}

\noindent
将一个范畴嵌入另一个范畴可能存在多种方式。有时它们是等价的，有时差异很大。一种方式可能将整个源范畴坍缩为单个对象，另一种方式可能将每个对象映射到不同对象，每个态射映射到不同态射。同一个蓝图可能以许多不同方式实现。自然变换(natural transformation)帮助我们比较这些实现方式。它们是函子之间的映射——保持函子性质的特殊映射。

考虑范畴$\cat{C}$和$\cat{D}$之间的两个函子$F$和$G$。若只关注$\cat{C}$中的单个对象$a$，它会被映射到两个对象：$F a$和$G a$。因此函子映射应当将$F a$映射到$G a$。

\begin{figure}[H]
  \centering
  \includegraphics[width=0.3\textwidth]{images/2_natcomp.jpg}
\end{figure}

\noindent
注意$F a$和$G a$都属于同一范畴$\cat{D}$中的对象。同一范畴内对象间的映射不应违背该范畴的结构特性。我们不想人为制造对象间的连接。因此使用现有连接（即态射(morphism)）是\emph{自然}的选择。自然变换就是一系列态射的选择：对每个对象$a$，选取从$F a$到$G a$的一个态射。若我们将自然变换记作$\alpha$，这个态射称为$\alpha$在$a$处的\newterm{分量}(component)，记作$\alpha_a$。

$$\alpha_a \Colon F a \to G a$$
请注意，这里的$a$是范畴$\cat{C}$中的对象，而$\alpha_a$是范畴$\cat{D}$中的态射。

如果对某些$a$，在$\cat{D}$中不存在从$F a$到$G a$的态射，则$F$和$G$之间不可能存在自然变换。

当然这只是故事的一半，因为函子不仅映射对象，还映射态射。那么自然变换如何处理这些映射？事实上态射的映射是固定的——在任意$F$和$G$之间的自然变换下，$F f$必须被转换为$G f$。更重要的是，两个函子对态射的映射极大地限制了定义兼容自然变换时的选择空间。考虑$\cat{C}$中两个对象$a$和$b$间的态射$f$，它会被映射到$\cat{D}$中的两个态射$F f$和$G f$：

\begin{gather*}
  F f \Colon F a \to F b \\
  G f \Colon G a \to G b
\end{gather*}
自然变换$\alpha$提供的两个额外态射完成了\emph{D}中的图示：

\begin{gather*}
  \alpha_a \Colon F a \to G a \\
  \alpha_b \Colon F b \to G b
\end{gather*}

\begin{figure}[H]
  \centering
  \includegraphics[width=0.4\textwidth]{images/3_naturality.jpg}
\end{figure}

\noindent
现在我们有了两条从$F a$到$G b$的路径。为保证两者相等，必须对任意$f$施加\newterm{自然性条件}(naturality condition)：

$$G f \circ \alpha_a = \alpha_b \circ F f$$
自然性条件是一个相当严格的约束。例如，若态射$F f$可逆，自然性就通过$\alpha_a$确定了$\alpha_b$。它会沿着$f$对$\alpha_a$进行\emph{传递}(transport)：

$$\alpha_b = (G f) \circ \alpha_a \circ (F f)^{-1}$$

\begin{figure}[H]
  \centering
  \includegraphics[width=0.4\textwidth]{images/4_transport.jpg}
\end{figure}

\noindent
若两个对象间存在多个可逆态射，则所有这些传递必须一致。不过一般情况下态射不可逆；但你可以看到两个函子间存在自然变换远非必然。因此由自然变换关联的函子数量稀少或丰富程度，可以揭示它们作用范畴的结构特征。当我们讨论极限(limit)和Yoneda引理(Yoneda lemma)时会看到具体例子。

从自然变换的分量角度看，可以说它将对象映射为态射。由于自然性条件的存在，也可以说它将态射映射为交换方图(commuting square)——每个$\cat{C}$中的态射对应$\cat{D}$中的一个交换自然方图。

\begin{figure}[H]
  \centering
  \includegraphics[width=0.4\textwidth]{images/naturality.jpg}
\end{figure}

\noindent
自然变换这一特性在大量范畴论构造中非常实用，这些构造通常包含交换图。通过精心选择函子，许多交换条件都可以转化为自然性条件。我们在讨论极限、余极限(colimit)和伴随(adjunction)时会看到这类例子。

最后，自然变换可用于定义函子的同构(isomorphism)。说两个函子自然同构几乎等同于说它们是相同的。\newterm{自然同构}(natural isomorphism)定义为所有分量都是同构（即可逆态射）的自然变换。

\section{多态函数}

我们讨论过函子（更确切地说是自函子(endofunctor)）在编程中的作用。它们对应将类型映射到类型的类型构造器。它们还将函数映射到函数，这种映射通过高阶函数\code{fmap}（或C++中的\code{transform}、\code{then}等）实现。

构造自然变换时，我们从一个对象开始，在这里是类型\code{a}。一个函子\code{F}将其映射到类型$F a$。另一个函子\code{G}将其映射到$G a$。自然变换\code{alpha}在\code{a}处的分量是从$F a$到$G a$的函数。用伪Haskell表示：

\begin{snipv}
alpha\textsubscript{a} :: F a -> G a
\end{snipv}
自然变换是一个对所有类型\code{a}定义的多态函数：

\src{snippet01}
在Haskell中\code{forall a}是可选的（实际上需要启用语言扩展\code{ExplicitForAll}）。通常你会这样写：

\src{snippet02}
请注意这实际上是参数化于\code{a}的一族函数。这是Haskell语法简洁性的又一体现。C++中类似的构造会稍显冗长：

\begin{snip}{cpp}
template<class A> G<A> alpha(F<A>);
\end{snip}
Haskell的多态函数与C++泛型函数还有一个更深刻的区别，体现在这些函数的实现和类型检查方式上。在Haskell中，多态函数必须对所有类型统一定义。一个公式必须适用于所有类型。这称为\newterm{参数多态}(parametric polymorphism)。

而C++默认支持\newterm{特设多态}(ad hoc polymorphism)，这意味着模板不必对所有类型都良好定义。模板是否适用于某个类型在实例化时决定——此时具体类型替代类型参数。类型检查被延迟，这往往导致难以理解的错误信息。

在C++中还有函数重载和模板特化机制，允许对不同类型定义同一函数的不同版本。Haskell通过类型类(type class)和类型族(type family)提供类似功能。

Haskell的参数多态有一个意外后果：任何形如：

\src{snippet03}
的多态函数（其中\code{F}和\code{G}是函子）自动满足自然性条件。用范畴论记号表示如下（$f$是函数$f \Colon a \to b$）：

$$G f \circ \alpha_a = \alpha_b \circ F f$$
在Haskell中，函子\code{G}对态射\code{f}的作用通过\code{fmap}实现。先用伪Haskell写出带显式类型注解的形式：

\begin{snipv}
fmap\textsubscript{G} f . alpha\textsubscript{a} = alpha\textsubscript{b} . fmap\textsubscript{F} f
\end{snipv}
由于类型推导，这些注解并非必需，以下等式成立：

\begin{snip}{text}
fmap f . alpha = alpha . fmap f
\end{snip}
这还不是真正的Haskell代码——函数相等性无法在代码中表达——但它是一个可供程序员进行等式推理或编译器进行优化的恒等式。

Haskell中自然性条件自动满足的原因与"免费定理"(free theorem)有关。用于在Haskell中定义自然变换的参数多态对实现有很强限制——一个公式适用于所有类型。这些限制转化为关于此类函数的等式定理。对于变换函子的函数，免费定理正是自然性条件。\footnote{
  你可以在我的博客中了解更多关于免费定理的信息：
  \href{https://bartoszmilewski.com/2014/09/22/parametricity-money-for-nothing-and-theorems-for-free/}{《Parametricity: Money for Nothing and Theorems for Free》}。}

前面提到过将Haskell中的函子视为广义容器的理解方式。我们可以延续这个类比，认为自然变换是将一个容器的内容重新打包到另一个容器的配方。我们不触碰元素本身：不修改它们，也不创造新元素。只是复制（某些）元素，有时多次复制，装入新容器。

自然性条件此时意味着：先通过\code{fmap}修改元素再重新打包，与先重新打包再通过新容器自己的\code{fmap}实现修改元素，二者等价。这两个操作——重新打包和\code{fmap}映射——是正交的。"一个移动鸡蛋，另一个煮熟它们。"

来看几个Haskell中的自然变换例子。第一个是列表函子与\code{Maybe}函子之间的变换。它返回列表的头部元素，但仅当列表非空时：

\src{snippet04}
这是一个参数多态于\code{a}的函数。它适用于任何类型\code{a}且无限制，是参数多态的典型例子。因此它是两个函子间的自然变换。但为了验证，我们来检查自然性条件。

\src{snippet05}
需要考虑两种情况；空列表：

\src{snippet06}

\src{snippet07}
和非空列表：

\src{snippet08}

\src{snippet09}
这里使用了列表的\code{fmap}实现：

\src{snippet10}
和\code{Maybe}的实现：

\src{snippet11}
一个有趣的情况是其中一个函子是平凡的\code{Const}函子。从或到\code{Const}函子的自然变换看起来就像在返回类型或参数类型上多态的函数。

例如，\code{length}可以视为从列表函子到\code{Const Int}函子的自然变换：

\src{snippet12}
这里\code{unConst}用于剥离\code{Const}构造器：

\src{snippet13}
当然实践中\code{length}定义为：

\src{snippet14}
这有效地隐藏了它作为自然变换的本质。

寻找从\code{Const}函子出发的参数多态函数稍难，因为它需要凭空创建值。我们最多能做到的是：

\src{snippet15}
另一个我们已见过的重要函子是\code{Reader}函子，它将在Yoneda引理中扮演关键角色。将其定义重写为\code{newtype}形式：

\src{snippet16}
它有两个类型参数，但仅对第二个参数协变函子：

\src{snippet17}
对每个类型\code{e}，可以定义一族从\code{Reader e}到任意其他函子\code{f}的自然变换。
后面我们会看到这些变换总是与\code{f e}的元素一一对应（参见\hyperref[the-yoneda-lemma]{Yoneda引理}）。

例如，考虑只有一个元素\code{()}的平凡单位类型\code{()}。\code{Reader ()}函子将任何类型\code{a}映射为函数类型\code{() -> a}。这些函数恰好对应从集合\code{a}中选取单个元素的所有可能选择。这样的函数数量与\code{a}中的元素数量相同。现在考虑从该函子到\code{Maybe}函子的自然变换：

\src{snippet18}
仅有两种这样的变换，\code{dumb}和\code{obvious}：

\src{snippet19}
和

\src{snippet20}
（你能对\code{g}做的唯一事情是将其应用于单位值\code{()}。）

确实，正如Yoneda引理预言的，它们对应着\code{Maybe ()}类型的两个元素\code{Nothing}和\code{Just ()}。我们之后会回到Yoneda引理——这只是个开胃小菜。

\section{超越自然性}

两个函子之间的参数化多态函数（包括\code{Const}函子这个边缘情况）总是自然变换。由于所有标准代数数据类型都是函子，这类类型之间的任意多态函数都是自然变换。

我们还可以使用函数类型，它们在其返回类型上具有函子性。我们可以用它们构建函子（如\code{Reader}函子），并定义高阶函数形式的自然变换。

然而，函数类型在其参数类型上并非协变的。它们是\newterm{逆变}的。当然，逆变函子与从对偶范畴来的协变函子是等价的。两个逆变函子之间的多态函数在范畴论意义上仍然是自然变换，只不过它们作用于从对偶范畴到Haskell类型的函子。

你可能还记得之前讨论过的逆变函子的例子：

\src{snippet21}
这个函子在\code{a}上是逆变的：

\src{snippet22}
我们可以编写一个从\code{Op Bool}到\code{Op String}的多态函数：

\src{snippet23}
但由于这两个函子不是协变的，这在$\Hask$范畴中并不是自然变换。然而，由于它们都是逆变的，它们满足"对偶"的自然性条件：

\src{snippet24}[b]
请注意这里的函数\code{f}必须与\code{fmap}使用的方向相反，这是由\code{contramap}的签名决定的：

\src{snippet25}
是否存在既不是协变也不是逆变的函子类型构造器？这里有一个例子：

\src{snippet26}
这不是函子，因为相同类型\code{a}同时出现在负变（逆变）和正变（协变）位置。你无法为此类型实现\code{fmap}或\code{contramap}。因此具有如下签名的函数：

\src{snippet27}
其中\code{f}是任意函子，不能构成自然变换。有趣的是，自然变换存在一种称为二自然变换(dinatural transformation)的广义形式，专门处理此类情形。当我们讨论端(ends)的概念时会重新探讨这个问题。

\section{函子范畴}

既然我们已经有了函子之间的映射——自然变换——很自然地会问：函子是否构成一个范畴？答案是肯定的！对于任意两个范畴$\cat{C}$和$\cat{D}$，都存在对应的函子范畴。该范畴中的对象是从$\cat{C}$到$\cat{D}$的所有函子，态射则是这些函子之间的自然变换。

我们需要定义两个自然变换的复合，这其实相当简单。自然变换的分量本身就是态射，而我们已经知道如何复合态射。

具体来说，取从函子$F$到$G$的自然变换$\alpha$，其在对象$a$处的分量是一个态射：
$$\alpha_a \Colon F a \to G a$$
现在我们要将$\alpha$与另一个从$G$到$H$的自然变换$\beta$复合。$\beta$在$a$处的分量是态射：
$$\beta_a \Colon G a \to H a$$
这两个态射可以复合，得到新的态射：
$$\beta_a \circ \alpha_a \Colon F a \to H a$$
我们将这个态射作为自然变换$\beta \cdot \alpha$在$a$处的分量，即$\beta$与$\alpha$按顺序复合的结果：
$$(\beta \cdot \alpha)_a = \beta_a \circ \alpha_a$$

\begin{figure}[H]
  \centering
  \includegraphics[width=0.4\textwidth]{images/5_vertical.jpg}
\end{figure}

\noindent
观察示意图可知，这种复合的结果确实是$F$到$H$的自然变换：
$$H f \circ (\beta \cdot \alpha)_a = (\beta \cdot \alpha)_b \circ F f$$

\begin{figure}[H]
  \centering
  \includegraphics[width=0.35\textwidth]{images/6_verticalnaturality.jpg}
\end{figure}

\noindent
自然变换的复合满足结合律，因为它们的分量（即普通态射）在其复合下具有结合性。

最后，对每个函子$F$，存在恒等自然变换$1_F$，其分量由恒等态射组成：
$$\id_{F a} \Colon F a \to F a$$
因此确实证明了函子构成一个范畴。

关于记号的说明：按照桑德斯·麦克兰恩的习惯，我使用点号表示上述自然变换的复合方式。问题在于自然变换有两种复合方式，这种称为垂直复合(vertical composition)，因为在相关示意图中函子通常纵向排列。垂直复合在定义函子范畴时至关重要。稍后我会解释水平复合(horizontal composition)。

\begin{figure}[H]
  \centering
  \includegraphics[width=0.3\textwidth]{images/6a_vertical.jpg}
\end{figure}

\noindent
范畴$\cat{C}$与$\cat{D}$之间的函子范畴写作$\cat{Fun(C, D)}$、$\cat{{[}C, D{]}}$，有时也写作$\cat{D^C}$。最后这种记号暗示函子范畴本身可能是某个范畴中的函数对象（指数对象）。这是否成立呢？

让我们回顾目前已构建的抽象层级。范畴论始于范畴本身，即由对象和态射组成的结构。范畴（严格来说是\emph{小范畴}，其对象构成集合）自身又是更高层范畴$\Cat$中的对象。该范畴的态射是函子。$\Cat$中的Hom集就是函子的集合，例如$\cat{Cat(C, D)}$表示范畴$\cat{C}$与$\cat{D}$之间所有函子的集合。

\begin{figure}[H]
  \centering
  \includegraphics[width=0.3\textwidth]{images/7_cathomset.jpg}
\end{figure}

\noindent
函子范畴$\cat{{[}C, D{]}}$同样是函子的集合（加上自然变换作为态射）。其对象与$\cat{Cat(C, D)}$的元素相同。更重要的是，作为范畴的函子范畴本身必然是$\Cat$的对象（恰好当$\cat{C}$和$\cat{D}$是小范畴时，其函子范畴也是小范畴）。我们在范畴的Hom集与其对象之间建立了对应关系，这正是上节讨论的指数对象的情形。现在来看如何在$\Cat$中构造这样的指数对象。

如前所述，要构造指数对象首先要定义乘积。在$\Cat$中这相对容易，因为小范畴本质上是对象的集合，而我们知道如何定义集合的笛卡尔积。因此乘积范畴$\cat{C\times D}$中的对象是形如$(c, d)$的有序对，其中$c$来自$\cat{C}$，$d$来自$\cat{D}$。类似地，这对象间的态射$(c, d)\to(c', d')$是态射对$(f, g)$，其中$f \Colon c \to c'$且$g \Colon d \to d'$。这些态射对按分量复合，且存在单位元即恒等态射对。简而言之，$\Cat$是一个完备的笛卡尔闭范畴，其中任意两个范畴都存在指数对象$\cat{D^C}$。这里的$\Cat$中的"对象"指范畴本身，因此$\cat{D^C}$就是一个范畴，我们可以将其等同于$\cat{C}$与$\cat{D}$之间的函子范畴。

\section{2-范畴}

解决了前面的问题之后，让我们更仔细地观察$\Cat$。根据定义，$\Cat$中的任意Hom集都是函子的集合。但是，正如我们所见，两个对象之间的函子具有比普通集合更丰富的结构。它们形成了一个范畴，其中自然变换作为态射。由于函子在$\Cat$中被视为态射，自然变换就是态射之间的态射。

这种更丰富的结构是$\cat{2}$-范畴的一个例子，这是范畴概念的一种推广：除了对象和态射（在此上下文中可能被称为1-态射）之外，还存在2-态射，即态射之间的态射。

当我们将$\Cat$视为$\cat{2}$-范畴时，具体表现为：

\begin{itemize}
  \tightlist
  \item
        对象：(小)范畴
  \item
        1-态射：范畴间的函子
  \item
        2-态射：函子间的自然变换
\end{itemize}

\begin{figure}[H]
  \centering
  \includegraphics[width=0.3\textwidth]{images/8_cat-2-cat.jpg}
\end{figure}

\noindent
对于范畴$\cat{C}$和$\cat{D}$之间的态射集，我们现在有了态射范畴——即函子范畴$\cat{D^C}$。我们有常规的函子复合：从$\cat{D^C}$取函子$F$与从$\cat{E^D}$取函子$G$复合得到$\cat{E^C}$中的$G \circ F$。但我们还在每个态射范畴内部有复合——即函子间自然变换（或2-态射）的垂直复合。

在$\cat{2}$-范畴中存在两种复合方式的情况下，自然产生一个问题：这两种复合如何相互作用？

让我们选取$\Cat$中的两个函子（即1-态射）：
\begin{gather*}
  F \Colon \cat{C} \to \cat{D} \\
  G \Colon \cat{D} \to \cat{E}
\end{gather*}
及其复合：
$$G \circ F \Colon \cat{C} \to \cat{E}$$
假设我们有两个分别作用于函子$F$和$G$的自然变换$\alpha$和$\beta$：
\begin{gather*}
  \alpha \Colon F \to F' \\
  \beta \Colon G \to G'
\end{gather*}

\begin{figure}[H]
  \centering
  \includegraphics[width=0.4\textwidth]{images/10_horizontal.jpg}
\end{figure}

\noindent
注意到我们无法对这对变换应用垂直复合，因为$\alpha$的靶对象与$\beta$的源对象不同。实际上它们属于两个不同的函子范畴：$\cat{D^C}$和$\cat{E^D}$。然而我们可以对函子$F'$和$G'$进行复合，因为$F'$的靶对象正是$G'$的源对象——即范畴$\cat{D}$。那么函子$G' \circ F'$与$G \circ F$之间有什么关系？

既然我们已经拥有$\alpha$和$\beta$，能否定义从$G \circ F$到$G' \circ F'$的自然变换？让我来勾勒构造过程。

\begin{figure}[H]
  \centering
  \includegraphics[width=0.5\textwidth]{images/9_horizontal.jpg}
\end{figure}

\noindent
按照惯例，我们从$\cat{C}$中的对象$a$开始。它的像分裂成$\cat{D}$中的两个对象：$F a$和$F'a$。还有一个连接这两个对象的态射（即$\alpha$的分量）：
$$\alpha_a \Colon F a \to F'a$$
当从$\cat{D}$到$\cat{E}$时，这两个对象进一步分裂为四个对象：$G (F a)$、$G'(F a)$、$G (F'a)$、$G'(F'a)$。我们还有构成正方形的四个态射。其中两个态射是自然变换$\beta$的分量：
\begin{gather*}
  \beta_{F a} \Colon G (F a) \to G'(F a) \\
  \beta_{F'a} \Colon G (F'a) \to G'(F'a)
\end{gather*}
另外两个则是$\alpha_a$在两个函子下的像（函子映射态射）：
\begin{gather*}
  G \alpha_a \Colon G (F a) \to G (F'a) \\
  G'\alpha_a \Colon G'(F a) \to G'(F'a)
\end{gather*}
这里有相当多的态射。我们的目标是寻找从$G (F a)$到$G'(F'a)$的态射，作为连接两个函子$G \circ F$和$G' \circ F'$的自然变换分量候选。事实上从$G (F a)$到$G'(F'a)$存在两条路径：
\begin{gather*}
  G'\alpha_a \circ \beta_{F a} \\
  \beta_{F'a} \circ G \alpha_a
\end{gather*}
幸运的是，这两条路径相等，因为我们构造的正方形恰好是$\beta$的自然性方块。

我们刚刚定义了从$G \circ F$到$G' \circ F'$的自然变换的一个分量。该变换的自然性证明相当直接，只要你有足够的耐心。

我们将这个自然变换称为$\alpha$和$\beta$的\newterm{水平组合}：
$$\beta \circ \alpha \Colon G \circ F \to G' \circ F'$$
再次说明，遵循麦克兰（Mac Lane）的记号，我用小圆圈表示水平组合，尽管你也可能看到星号被用于相同目的。

这里有个范畴论的经验法则：每当你遇到复合操作时，都应该寻找其所属的范畴。我们已有自然变换的垂直组合，它构成了函子范畴的一部分。但水平组合呢？它存在于哪个范畴中？

解决这个问题的方法是横向观察$\Cat$。不要将自然变换视为函子间的箭头，而是将其视为范畴间的箭头。自然变换处于两个范畴之间，而这两个范畴通过它变换的函子相连。我们可以认为它连接着这两个范畴。

\begin{figure}[H]
  \centering
  \includegraphics[width=0.5\textwidth]{images/sideways.jpg}
\end{figure}

\noindent
让我们聚焦$\Cat$的两个对象——范畴$\cat{C}$和$\cat{D}$。存在一组自然变换，在连接$\cat{C}$到$\cat{D}$的函子之间。这些自然变换就是我们新的从$\cat{C}$到$\cat{D}$的箭头。同理，存在在连接$\cat{D}$到$\cat{E}$的函子之间的自然变换，我们可以将其视为从$\cat{D}$到$\cat{E}$的新箭头。水平组合就是这些箭头的复合。

我们还有从$\cat{C}$到$\cat{C}$的恒等箭头。它是将$\cat{C}$上的恒等函子映射到自身的恒等自然变换。注意水平组合的恒等元同时也是垂直组合的恒等元，但反过来则不成立。

最后，这两种组合满足交换律：
$$(\beta' \cdot \alpha') \circ (\beta \cdot \alpha) = (\beta' \circ \beta) \cdot (\alpha' \circ \alpha)$$
在此引用桑德斯·麦克兰（Saunders Mac Lane）的话：读者或许会乐于画出证明此结论所需的明显图示。

未来可能会用到的另一个记法值得一提。在这种横向解释的$\Cat$中，有两种方式可以在对象间移动：使用函子或使用自然变换。然而，我们可以重新将函子箭头解释为某种特殊的自然变换：作用于该函子上的恒等自然变换。因此你经常会看到如下记法：
$$F \circ \alpha$$
其中$F$是从$\cat{D}$到$\cat{E}$的函子，$\alpha$是两个从$\cat{C}$到$\cat{D}$的函子间的自然变换。由于你无法将函子与自然变换直接复合，这应被解读为$\alpha$后接恒等自然变换$1_F$的水平组合。

同理：
$$\alpha \circ F$$
表示$1_F$后接$\alpha$的水平组合。

\section{结论}

本书的第一部分到此结束。我们已经学习了范畴论的基本术语。你可以将对象和范畴视为名词；而将态射、函子和自然变换视为动词。态射连接对象，函子连接范畴，自然变换连接函子。

但我们还注意到，在某个抽象层级上的操作行为，在更高一层级会表现为对象。一组态射可转化为函数对象。作为对象后，它又能成为其他态射的源或目标。这就是高阶函数的核心思想。

函子将对象映射到对象上，因此我们可以将其用作类型构造器或参数化类型。函子同时也映射态射，因此它是高阶函数——\code{fmap}。存在一些简单的函子，如 \code{Const}、乘积和余乘积，它们可用于生成多种代数数据类型。函数类型也具有函子性，既协变又逆变，可用于扩展代数数据类型。

函子可被视为函子范畴中的对象。作为对象后，它们又成为态射（自然变换）的源和目标。自然变换是一种特殊的多态函数。

\section{挑战}

\begin{enumerate}
  \tightlist
  \item
        定义从\code{Maybe}函子到列表函子的一个自然变换。证明其满足自然条件。
  \item
        在\code{Reader ()}与列表函子之间定义至少两个不同的自然变换。有多少种不同的\code{()}列表？
  \item
        用\code{Reader Bool}和\code{Maybe}继续前一个练习。
  \item
        证明自然变换的水平合成满足自然条件（提示：使用分量）。这是练习图示追踪的好机会。
  \item
        写一篇短文阐述你如何享受绘制证明交换律所需的直观图表过程。
  \item
        为不同\code{Op}函子之间变换的对偶自然条件创建几个测试用例。以下是一个示例：

        \begin{snip}{haskell}
op :: Op Bool Int
op = Op (\x -> x > 0)
\end{snip}
        以及

        \begin{snip}{haskell}
f :: String -> Int
f x = read x
\end{snip}
\end{enumerate}